% Research Note of the AAS (RNAAS) draft.
% Every number in this file comes from \input{numbers} / \input{table}, which
% research/figures/rnaas_note.py generates from research/results/phase1_full/.
% Build: research/papers/rnaas/build.sh  (copies the generated files, runs pdflatex)
\documentclass{aastex701}

\input{numbers}

\begin{document}

\title{Keeping a ``None of the Above'' Hypothesis Matters More Than Planning With It: A Controlled Test for Transient Follow-Up}

\author[0000-0000-0000-0000]{Joyce Malik}
\affiliation{[Department], [Institution], [City], India}
\email{corresponding.author@example.org}
% TODO before submission: real ORCID, affiliation, email, and any co-authors.

\begin{abstract}
Information-driven follow-up planners for astronomical transients usually assume the event belongs to one of a fixed set of classes. We test, in a controlled toy problem, what happens when it does not. Across \NEvents{} simulated fading transients (\NSeeds{} seeds), a closed-world expected-information-gain planner was confidently wrong on \BCwWith\% of events from a withheld class. Adding an explicit unmodeled hypothesis $h_0$, represented by a Gaussian process, to the posterior removes most of this failure. Planning with $h_0$ helped only modestly: the open-world planner flagged \PFlagWith\% of withheld events, against \BHFlagWith\% for the closed-world planner's measurements rescored with $h_0$. Randomly timed measurements rescored the same way flagged \RndHFlagWith\%. Information-gain planning, even with $h_0$, chose measurements that exposed model inadequacy less well than random timing, so inadequacy detection has to be designed into the planning objective explicitly.
\end{abstract}

\section{Motivation}
Wide-field surveys \citep{Bellm2019,Ivezic2019} discover more transients than follow-up resources can observe, so which measurement to take next is a planning problem. Expected information gain (EIG) about the class label \citep{Lindley1956,Chaloner1995} is a natural objective, but it presumes the true class is among those modeled. Rare and new phenomena are precisely the ones that are not. We isolate this assumption in a deliberately small problem.

\section{Setup}
Each event is drawn uniformly from three classes: two \emph{modeled} classes with flux $A\,e^{-t/\tau}$ ($\tau = 1$ and $4$\,d), and one \emph{withheld} class, a Gaussian rise peaking at 1.5\,d followed by exponential decay, which no planner's model contains. Amplitudes are log-uniform on $[5, 20]$ with Gaussian noise $\sigma = 1$. Each strategy takes 5 measurements chosen from 20 log-spaced epochs between 0.2 and 10\,d. Noise is pre-drawn per epoch, so all strategies see identical events and noise.

The closed-world belief is an exact grid posterior over class and amplitude (25 log-spaced values). The open-world posterior adds $h_0$, scored by the marginal likelihood of a zero-mean Gaussian process (squared-exponential kernel, length scale 1.2\,d, variance 100; fixed, not fitted), with a uniform prior over \{modeled classes, $h_0$\}. EIG is computed by exact quadrature. Three strategies choose epochs: random; closed-world EIG; and open-world EIG, whose predictive mixture includes the $h_0$ Gaussian process. Because every run records both posteriors at every step, we also score the random and closed-world EIG measurement sequences with the open-world posterior after the fact. This separates the effect of \emph{having} $h_0$ from that of \emph{planning} with it. A decision is ``confident'' when its posterior exceeds 0.9. Code and per-event outputs are public (see below).

\section{Results}
Table~\ref{tab:res} gives pooled results with 95\% Wilson intervals over \NModeled{} modeled-class and \NWithheld{} withheld-class events. Modeled-class accuracy was between \RndHAccMod\% and \BAccMod\% for all strategies. On withheld events, every closed-world posterior must name a modeled class, and closed-world EIG did so with high confidence on \BCwWith\% of them. With the open-world posterior, confident errors fell to \BHCwWith\% for the same closed-world EIG measurements, \RndHCwWith\% for random measurements and \PCwWith\% for the open-world planner. On paired withheld events, the open-world planner flagged $h_0$ where post-hoc closed-world EIG did not in \PairPOnly{} cases, and the reverse occurred in \PairBOnly{} (exact McNemar $p = \PairP$). The ordering of $h_0$-flag rates (random above open-world EIG above closed-world EIG) held in each of the three seeds. The area under the ROC curve of $P(h_0)$ separating withheld from modeled events was \RndHAuroc{} (random), \BHAuroc{} (closed-world EIG) and \PAuroc{} (open-world EIG).

\begin{table}[h]
\centering
\caption{Pooled outcomes (percent, 95\% Wilson interval). ``Post hoc'': measurements chosen without $h_0$, scored with the open-world posterior.}
\label{tab:res}
\scriptsize
\input{table}
\end{table}

\section{Interpretation and limitations}
In this setting, most of the benefit of an open-world treatment comes from the \emph{posterior}. Planning with $h_0$ in the predictive mixture adds a small, significant improvement over closed-world planning, but both information-gain planners detected inadequacy less often, and separated it less well, than random timing. We have not isolated why. A plausible reading is that EIG favours epochs that discriminate among hypotheses, which is not the same as epochs that expose misfit. The practical implication is that follow-up planners should carry an explicit adequacy term in their utility, such as the mutual information between a measurement and a misspecification indicator, rather than relying on class-level information gain.

This is a toy problem: two exponential classes, one withheld shape, fixed noise, fixed Gaussian process hyperparameters, a uniform class prior, no observability constraints and a single threshold. The results establish a failure mode and a design requirement, not performance in realistic follow-up. A study with published light-curve templates, real orbital visibility and a preregistered evaluation is in progress.

\section*{Code availability}
Simulator, planners, raw outputs and the script that generates every number here: \url{https://github.com/joycemalik/orbitalnet-dashboard} (branch \texttt{research/open-world-planner}). % TODO: replace with the Zenodo DOI of the tagged release.

\begin{thebibliography}{}
\bibitem[Bellm et al.(2019)]{Bellm2019} Bellm, E.~C., et al. 2019, PASP, 131, 018002
\bibitem[Chaloner \& Verdinelli(1995)]{Chaloner1995} Chaloner, K., \& Verdinelli, I. 1995, StaSc, 10, 273
\bibitem[Ivezi\'c et al.(2019)]{Ivezic2019} Ivezi\'c, \v{Z}., et al. 2019, ApJ, 873, 111
\bibitem[Lindley(1956)]{Lindley1956} Lindley, D.~V. 1956, AnMS, 27, 986
\end{thebibliography}

\end{document}
